\appendixitem{DOMINANT POLYNOMIAL MAPS}

Let V, W be finite dimensional vector spaces. Let f be a polynomial map from V to W (Algebra, Chap. IV, §5, no. 10, Def. 4). If \(\psi \in A_{W}\) , \(\psi \circ f \in A_{V}\) (loc. cit., Prop. 17). The map \(\psi \mapsto \psi \circ f\) is a homomorphism from \(A_{W}\) to \(A_{V}\) , said to be associated to f.~Its kernel consists of the functions \(\psi \in A_{W}\) which vanish on \(f(V)\) (and hence also on the closure of \(f(V)\) in the Zariski topology).

DEFINITION 1. A polynomial map \(f: \mathrm{V} \to \mathrm{W}\) is said to be dominant if the homomorphism from \(\mathrm{A_W}\) to \(\mathrm{A_V}\) associated to \(f\) is injective.

In view of the preceding, \(f\) is dominant if and only if \(f(\mathrm{V})\) is dense in W in the Zariski topology.

PROPOSITION 3. Assume that k is algebraically closed. Let \(f : V \to W\) be a dominant polynomial map. The image under f of any dense open subset of V contains a dense open subset of W.

It suffices to prove that, for every non-zero element \(\varphi\) of \(A_{V}\) , \(f(V_{\varphi})\) contains a dense open subset of W. Identify \(A_{W}\) with a subalgebra of \(A_{V}\) by means of the homomorphism associated to \(f\) . There exists a non-zero element \(\psi\) of \(\mathrm{A_W}\) such that every homomorphism \(w: \mathrm{A_W} \to k\) which does not annihilate \(\psi\) extends to a homomorphism \(v: \mathrm{A_V} \to k\) which does not annihilate \(\varphi\) (Commutative Algebra, Chap. V, §3, no. 1, Cor. 3 of Th. 1). Now such a \(w\) (resp. \(v\) ) can be identified with an element of \(\mathrm{W}_{\psi}\) (resp. of \(\mathrm{V}_{\varphi}\) ) and to say that \(v\) extends \(w\) means that \(f(v) = w\) . Hence, \(\mathrm{W}_{\psi} \subset f(\mathrm{V}_{\varphi})\) . Q.E.D.

Let \(f: \mathrm{V} \to \mathrm{W}\) be a polynomial map, and \(x_0 \in \mathrm{V}\) . The map \(h \mapsto f(x_0 + h)\) from \(\mathrm{V}\) to \(\mathrm{W}\) is polynomial. Decompose it into a finite sum of homogeneous polynomial maps:

\[
f (x _ {0} + h) = f (x _ {0}) + \mathrm{D} _ {1} (h) + \mathrm{D} _ {2} (h) + \dots
\]

where \(D_{i}: V \to W\) is homogeneous of degree i (Algebra, Chap. IV, §5, no. 10, Prop. 19). The linear map \(D_{1}\) is called the tangent linear map of f at \(x_{0}\) . We denote it by \(\mathrm{D}f(x_{0})\) .

PROPOSITION 4. Let \(f: \mathrm{V} \to \mathrm{W}\) be a polynomial map. Assume that there exists \(x_0 \in \mathrm{V}\) such that \((\mathrm{D}f)(x_0)\) is surjective. Then \(f\) is dominant.

Applying a translation in V and one in W, we can assume that \(x_{0} = 0\) and \(f(x_{0}) = 0\) . The decomposition of f as a sum of homogeneous elements can then be written

\[
f = f _ {1} + f _ {2} + \dots \text {with} \deg f _ {i} = i,
\]

and the linear map \(f_{1}\) is surjective by hypothesis. Suppose that f is not dominant. Then there exists a non-zero element \(\psi\) of \(A_{W}\) such that \(\psi \circ f = 0\) . Let \(\psi = \psi_{m} + \psi_{m+1} + \cdots\) be the decomposition of \(\psi\) into homogeneous elements, with \(\deg \psi_{i} = i\) and \(\psi_{m} \neq 0\) . Then

\[
\begin{array}{l} 0 = \psi \circ f = \psi_ {m} \circ f + \psi_ {m + 1} \circ f + \dots \\ = \psi_ {m} \circ f _ {1} + \rho , \end{array}
\]

where \(\rho\) is a sum of homogeneous polynomial maps of degrees \textgreater{} m. It follows that \(\psi_{m} \circ f_{1} = 0\) . Since \(f_{1}\) is surjective, \(\psi_{m} = 0\) , a contradiction.

COROLLARY. If k is algebraically closed and if f satisfies the assumptions of Prop. 4, the image under f of any dense open subset of V contains a dense open subset of W.

This follows from Props. 3 and 4.
% semantic-fragments: legacy-input-seam
\relax{}
